%\chapter{ABSTRAK}
\clearpage
\phantomsection
\addcontentsline{toc}{chapter}{INTISARI}
\begin{center}
   \textbf{\large{\judulid}}\\[0.5cm]
   Oleh :\\
   \penulis\\
   \nim\\[2em]
    \textbf{INTISARI}\\[0.5cm]
\end{center}

Hotel kapsul merupakan salah satu bentuk akomodasi modern yang tengah populer, terutama di kawasan perkotaan dengan tingkat mobilitas tinggi. Konsepnya yang mengedepankan efisiensi ruang, biaya terjangkau, serta pengalaman menginap yang unik menjadikan hotel kapsul pilihan menarik bagi wisatawan domestik maupun mancanegara. Namun, efisiensi fisik tersebut perlu didukung oleh sistem manajemen operasional yang setara dari sisi teknologi agar pengelolaan hotel berjalan optimal. Penelitian ini bertujuan untuk merancang dan mengembangkan aplikasi manajemen tiket hotel kapsul berbasis web yang terintegrasi dan responsif terhadap kebutuhan operasional di era digital. Aplikasi ini dikembangkan dengan mengedepankan kemudahan proses pemesanan, pembayaran, dan pengelolaan tiket baik dari sisi penyedia layanan maupun tamu hotel. Sistem ini mendukung tiga jenis pengguna, yaitu superadmin sebagai pemilik waralaba (Tamu.id), admin sebagai pengelola masing-masing hotel kapsul, dan pengguna sebagai tamu. Fitur utama yang diimplementasikan meliputi autentikasi pengguna (termasuk login tanpa kata sandi melalui integrasi Google OAuth2), pemesanan tiket dengan batas waktu dan mekanisme optimistic locking untuk mencegah race condition, pembayaran daring secara real-time melalui QRIS yang diintegrasikan dengan webhook Midtrans via WebSocket, serta sistem kunci kamar berbasis QR code dengan metode regenerasi berantai (chained QR) yang memungkinkan integrasi dengan perangkat IoT secara offline ready. Selain itu, sistem ini juga mendukung pemberian voucher diskon dengan sistem aturan majemuk (multiple ruling), skema referral untuk promosi, pengaturan harga dinamis berdasarkan rentang waktu tertentu (temporal pricing), sinkronisasi data pemesanan dari online travel agency (OTA) eksternal, dan fitur manajemen konten (CMS) untuk kebutuhan informasi hotel. Pengembangan sistem ini menggunakan pendekatan metodologi Agile dengan model Kanban untuk menjaga fleksibilitas dan ketepatan waktu pengembangan. Hasil akhir dari penelitian ini adalah sebuah aplikasi manajemen tiket hotel kapsul berbasis web yang mampu meningkatkan efisiensi operasional hotel, memperkaya pengalaman pengguna, dan menyediakan solusi teknologi yang siap diintegrasikan ke dalam ekosistem digital perhotelan modern.

\noindent Kata kunci: \katakunci